\section{Introduction}

Post-training has evolved from a simple one-pass stage in LLM production into a widely used infrastructure workload. As LLMs move toward trillion-scale parameter counts, real-world deployment, and lifelong learning from experience~\citep{yao_second_half_2025,silver_sutton_experience_2025}, frontier model developers increasingly emphasize the complexity of building reliable frameworks for training modern agentic LLM capabilities~\citep{deepseek_v4_release_2026,deepseekv4_2026,glm5_2026,kimi_k25_2026,minimax_m27_2026,qwen35_2026,openai_gpt55_2026,anthropic_opus45_2025}. These workloads introduce a range of infrastructure challenges, including hardware resource management, distributed-computing scheduling, and version control for model artifacts. Traditional infrastructures rely on copying or serving a full fine-tuned checkpoint for each model variant are increasingly difficult to scale under the modern demands of continuous training, deployment and agentic reinforcement learning.

% \clearpage

\begin{figure}[H]
    \centering
    \resizebox{\textwidth}{!}{\begin{tikzpicture}[x=1cm,y=1cm,>=stealth,font=\sffamily]
  \tikzstyle{panel}=[draw=mindlabfg, very thick, rounded corners=1.4mm, fill=white]
  \tikzstyle{mintpanel}=[draw=mindlabfg, very thick, rounded corners=1.4mm, fill=mindlabblue!12]

  \tikzstyle{box}=[
    draw=mindlabfg,
    thick,
    rounded corners=0.9mm,
    fill=white,
    align=center,
    inner sep=3pt
  ]

  \tikzstyle{softbox}=[
    box,
    fill=mindlabbluepale!45,
    minimum width=2.34cm,
    minimum height=0.62cm,
    text width=2.10cm
  ]

  \tikzstyle{groupbox}=[
    draw=mindlabfg,
    thick,
    rounded corners=0.9mm,
    fill=mindlabbluepale!45,
    align=center,
    inner sep=3pt,
    font=\scriptsize\sffamily
  ]

  \tikzstyle{adapter}=[
    draw=mindlabfg,
    thick,
    rounded corners=0.8mm,
    fill=mindlabblue!22,
    minimum width=0.58cm,
    minimum height=0.24cm,
    inner sep=0.6pt,
    font=\scriptsize\sffamily,
    text=mindlabfg
  ]

  \tikzstyle{base}=[
    draw=mindlabfg,
    thick,
    rounded corners=0.9mm,
    fill=mindlabfg!8,
    align=center,
    inner sep=3pt,
    text width=2.55cm
  ]

  \tikzstyle{title}=[
    font=\bfseries\small\sffamily,
    align=center,
    text=mindlabfg
  ]

  \tikzstyle{bigtag}=[
    font=\bfseries\large\sffamily,
    align=center,
    text=mindlabfg
  ]

  \tikzstyle{label}=[
    font=\scriptsize\sffamily,
    align=center,
    text=mindlabfg,
    fill=white,
    inner sep=1pt
  ]

  \tikzstyle{plainlabel}=[
    font=\scriptsize\sffamily,
    align=center,
    text=mindlabfg,
    inner sep=1pt
  ]

  \tikzstyle{arrow}=[->, very thick, draw=mindlabfg]

  \path[use as bounding box] (0.10,-1.02) rectangle (14.10,7.95);

  % User intent
  \node[panel, minimum width=12.85cm, minimum height=1.76cm] (user) at (7.10,7.00) {};
  \node[title] at (7.10,7.65) {User Intent};

  \foreach \x/\main/\sub in {
    2.40/Base Model/{Qwen, DeepSeek},
    5.53/Data + Reward/{tasks, tools, envs},
    8.67/LoRA RL Recipe/{rank, modules, loss},
    11.80/Evaluate + Serve/{scores, rollout, serve}
  } {
    \node[
      box,
      minimum width=2.82cm,
      minimum height=0.62cm,
      text width=2.50cm
    ] at (\x,6.82)
      {{\bfseries\footnotesize \main}\\[-1pt]
       {\scriptsize\color{mindlabfg!62}\sub}};
  }

  % MinT
  \node[mintpanel, minimum width=12.85cm, minimum height=1.42cm] (mint) at (7.10,5.06) {};
  \node[bigtag] at (1.55,5.06) {MinT};

  \foreach \x/\txt in {
    3.82/API,
    6.52/Queue,
    9.22/Policy State,
    11.92/Revisions
  } {
    \node[softbox] at (\x,5.06)
      {{\footnotesize \txt}};
  }

  % Policy population
  \node[panel, minimum width=12.85cm, minimum height=2.18cm] (population) at (7.10,2.91) {};
  \node[title] at (7.10,3.73) {Policy Population over Shared Base Deployments};

  \node[base, minimum width=2.95cm, minimum height=0.82cm] (bases) at (4.29,2.78)
    {{\footnotesize Resident Base}\\[-1pt]
     {\footnotesize Deployments}};

  \node[plainlabel, font=\bfseries\large\sffamily] at (6.23,2.78) {$+$};

  \node[adapter] at (7.29,2.96) {$r_1$};
  \node[adapter] at (8.05,2.96) {$r_2$};
  \node[adapter] at (8.81,2.96) {$r_3$};
  \node[adapter] at (9.57,2.96) {$r_4$};
  \node[adapter] at (10.33,2.96) {$r_5$};
  \node[adapter] at (11.09,2.96) {$\cdots$};

  \node[plainlabel] at (9.19,2.44) {train / evaluate / serve / rollback};

  % Infrastructure
  \node[panel, minimum width=12.85cm, minimum height=2.26cm] (infra) at (7.10,0.34) {};
  \node[title] at (7.10,1.20) {Infrastructure Complexity Hidden Underneath};

  \node[groupbox, minimum width=2.62cm, minimum height=1.10cm, text width=2.44cm] at (2.46,0.15)
    {{\bfseries\scriptsize Scheduler}\\[-1pt]
     {\scriptsize\color{mindlabfg!62}queue / admission}\\[-1pt]
     {\scriptsize\color{mindlabfg!62}placement}};

  \node[groupbox, minimum width=2.62cm, minimum height=1.10cm, text width=2.44cm] at (5.55,0.15)
    {{\bfseries\scriptsize Fault Tolerance}\\[-1pt]
     {\scriptsize\color{mindlabfg!62}records / op state}\\[-1pt]
     {\scriptsize\color{mindlabfg!62}retries}};

  \node[groupbox, minimum width=2.62cm, minimum height=1.10cm, text width=2.44cm] at (8.65,0.15)
    {{\bfseries\scriptsize Adapter Lifecycle}\\[-1pt]
     {\scriptsize\color{mindlabfg!62}restore / export}\\[-1pt]
     {\scriptsize\color{mindlabfg!62}revision visibility}};

  \node[groupbox, minimum width=2.62cm, minimum height=1.10cm, text width=2.44cm] at (11.74,0.15)
    {{\bfseries\scriptsize Serving Residency}\\[-1pt]
     {\scriptsize\color{mindlabfg!62}storage / cache}\\[-1pt]
     {\scriptsize\color{mindlabfg!62}GPU slots}};

  \draw[arrow] (user.south) -- (mint.north);
  \draw[arrow] (mint.south) -- (population.north);
  \draw[arrow] (population.south) -- (infra.north);
\end{tikzpicture}}
    \caption{MinT overview. User intent selects a base model, data and rewards, a LoRA RL recipe, and evaluation or serving targets. MinT turns those inputs into queued work, policy records, and exported revisions, then manages a policy population over shared base deployments while scheduler, fault-tolerance, adapter-lifecycle, and serving-residency mechanisms remain behind the service interface.}
    \label{fig:mint_overview}
\end{figure}

To address these challenges, we propose the MindLab Toolkit (MinT), which uses Low-Rank Adaptation (LoRA) adapters as the basic policy units for post-training and online serving. In MinT, the base model remains resident, while task behaviors, product branches, experimental versions, tenant-specific variants, and rollback points are represented by different adapters~\citep{lora2022,lora_without_regret_2025}. As shown in \Cref{fig:mint_overview}, MinT transforms user intent into queued jobs, policy records, and exported adapter revisions on top of shared base-model deployments. Instead of repeatedly moving, loading, or copying full models between training and serving, the system transfers and manages only compact LoRA parameters. Following the service-interface practice of Tinker~\citep{tinker2025,tinker_cookbook}, MinT connects rollout, update, export, evaluation, serving, and rollback into a LoRA-centered workflow. It hides distributed training, model sharding, tensor-layout conversion, adapter admission, and serving access behind a simple service interface, making large-scale LoRA-based reinforcement learning easier to run, reproduce, and deploy.

This adapter-centered design changes what crosses the training-serving boundary. Full fine-tuning moves a full checkpoint for each trained variant. Merge-based LoRA reduces training memory, but still folds the adapter back into the base model and moves a merged checkpoint before inference. MinT instead exports the updated LoRA as a serving-compatible adapter revision, checks that it matches the resident base model, and loads it into an inference engine that already holds that base. \Cref{fig:mint_handoff_paths} illustrates this difference: MinT moves adapter revisions, not full model checkpoints.

\begin{figure}[H]
    \centering
    \resizebox{\textwidth}{!}{\input{figures/mint_teaser}}
    \caption{What crosses from training to serving under three adaptation paths. Blue marks the trained parameters. Full fine-tuning moves a full checkpoint for each variant. Merge-based LoRA trains adapters beside a resident base, then moves merged full checkpoints. MinT trains adapters beside a resident base and moves only adapter revisions to an inference engine that already holds the compatible base.}
    \label{fig:mint_handoff_paths}
\end{figure}

In this report, we introduce MinT's capabilities follow three scaling axes. \textit{Scale Up} supports LoRA RL on frontier-scale dense and MoE architectures, with model-parallel training and serving paths validated beyond 1T total parameters~\citep{qwen3_2025,deepseekv4_2026,glm5_2026,kimi_k2_2025}. This axis keeps the same adapter-revision path usable when the base requires tensor-parallel or expert-parallel placement, sparse-route consistency, and distributed adapter export. \textit{Scale Down} minimizes the training-serving handoff by moving only the exported LoRA adapter, which can be less than 1\% of the base-model size in compact rank-1 settings, eliminating full-checkpoint materialization entirely. In our measurements, adapter-only handoff reduces the measured handoff step by $18.3\times$ on a 4B dense model and by $2.85\times$ on a 30B MoE model; under the same resident-base allocation, concurrent multi-policy GRPO shortens wall time by $1.77\times$ and $1.45\times$, respectively, without increasing peak memory. \textit{Scale Out} expands the policy namespace while keeping engine-local execution bounded. MinT separates durable policy addressability from CPU/GPU hot working sets, builds on modern multi-adapter serving support~\citep{vllm2023,punica2024,slora2023,dlora2024,loraserve2025}, and treats first-touch adapter loading as scheduled service work. A tensor-parallel serving deployment supports $10^6$-scale addressable policy catalogs, with measured single-engine sweeps through 100K entries and thousand-adapter active waves at cluster scale; packed MoE LoRA tensors improve live engine loading by $8.5$--$8.7\times$ by reducing small-object fanout on the cold-load path. Together, these axes make continuously evolving LoRA post-training cheaper to operate, easier to reproduce, and better matched to multi-tenant policy populations.

MinT contributes four concrete pieces:

\begin{itemize}
    \item \textbf{Adapter lifecycle from update to serving.} MinT exports trained LoRAs as PEFT adapter revisions, including sharded-to-serving conversion, compatibility checks, rollout records, sampler loading, served-result attribution, and rollback. The adapter revision becomes the fixed behavior selected by rollout, evaluation, online serving, and recovery.

    \item \textbf{Large-scale multi-LoRA RL training.} MinT time-slices LoRA training sessions over resident shared bases and supports both single-worker PEFT and distributed Megatron training paths. It reports end-to-end LoRA RL on large dense and MoE deployments, including 235B-A22B-scale training and a 1T-class MoE countdown-task path.

    \item \textbf{Policy-population multi-LoRA serving.} MinT serves exported adapters through shared-base vLLM engines, separates durable policy addressability from CPU/GPU hot working sets, and treats cold loading as scheduled work for cache misses. Its packed MoE LoRA serving representation reduces small-object fanout and speeds live engine loading by 8.5--8.7$\times$.

    \item \textbf{Public reproducibility paths.} MinT provides a Tinker-compatible API~\citep{tinker2025,tinker_cookbook} and uses mint-cookbook recipes~\citep{mint_cookbook2026} to reproduce SFT, preference optimization, rollout-based RL, and AutoResearch examples through the same adapter lifecycle.
\end{itemize}

%%%


% 随着大语言模型进入更广泛的产品场景，用户对模型能力的需求持续上升，推动模型参数规模、后训练任务和生产部署流程不断扩张。模型规模的增长不仅意味着更高的算力消耗，也显著提高了训练与服务系统的整体复杂度：在硬件层面，系统需要管理大规模 GPU 集群，以及机房供电、散热、网络和资源调度等基础设施；在训练软件层面，系统需要协调分布式训练、任务编排、容错恢复、数据传输和推理服务；在模型工程层面，团队还需要维护模型版本、数据版本、实验分支、评测结果和回滚路径。尤其是在长程学习、工具使用和 agentic post-training 等场景中，训练、推理、环境交互和服务反馈被更紧密地耦合在一起，版本切换操作变得更加频繁和昂贵，使传统面向单个最终 checkpoint 的基础设施越来越难以支撑持续迭代。因此，LLM 后训练亟需一种更轻量、更可管理、也更用户友好的基础设施抽象。在实践中，LoRA 由于能够在冻结基础模型的同时，用小规模参数承载新的模型行为，被广泛认为是缓解这一复杂度的一条有潜力的路径。 



% MinT 正是围绕这一思路设计的基础设施系统。它并不把 LoRA 仅仅视为一种节省显存的训练技巧，而是将其提升为模型后训练中的基本策略单元：基础模型保持常驻，而具体的训练行为、产品分支、实验版本和回滚点则由不同的 LoRA adapter 承载。这样一来，系统在训练和服务之间不再需要频繁移动、加载或复制完整模型，而只需要传输和管理小规模的 LoRA 参数。进一步的，通过follow the practice of Tinker framework, MinT 将 rollout、参数更新、推理服务、版本导出和回滚等流程统一到一个以 LoRA 为中心的工作流中，并将底层复杂的分布式训练、模型切分、张量布局转换和服务接入隐藏在简洁的服务接口之后。通过这种方式，MinT 将训练更新和服务切换的操作工作在 LoRA adapter 上，避免频繁处理完整模型，从而降低多分支、多版本和多租户训练服务的系统复杂度。



% MinT 的能力可以从三个扩展维度来理解。

% 首先，**Scale Up** 关注模型规模的向上扩展，使 LoRA-based RL 能够运行在更大的 dense 和 MoE 架构之上，支持 Qwen、DeepSeek-style 等大模型家族，并在 GLM-5.1、Kimi K2 等前沿训练路径上进行验证。

% 其次，**Scale Down** 关注训练与服务之间的开销压缩。由于基础模型保持常驻，系统只需要在训练器和推理引擎之间传输 LoRA adapter，而不需要反复物化、搬移和加载完整 checkpoint。在极端情况下，rank-1 LoRA 的规模可以小到基础模型的约 1%，从而显著降低模型更新和服务切换成本。

% 最后，**Scale Out** 关注横向并发扩展。MinT 通过多级 adapter cache 同时管理和服务大量 LoRA 策略版本，使同一个基础模型之上可以承载多任务、多租户、多实验分支，乃至百万级 policy revisions。

% 通过这三个维度，MinT 不仅降低了单次训练或单次服务的成本，更将大规模、多版本、持续迭代的 LoRA 后训练转化为一种可管理、可复现、可扩展的基础设施能力。

% 基于上述设计，MinT 提供了三类核心系统能力。

% 第一，MinT 使超大模型上的 LoRA 强化学习训练更加可行。在大模型后训练中，多个实验分支往往需要共享同一个基础模型，但又各自维护不同的策略更新。MinT 通过让基础模型常驻，并在其上切换不同 LoRA adapter，使多个训练会话能够复用同一套底层模型资源。系统已验证其能够支持 Qwen3-235B-A22B 等大规模模型上的端到端 LoRA RL，并通过 Kimi K2 1T MoE 等更大规模部署路径进一步展示了这种架构的扩展潜力。

% 第二，MinT 降低了多版本 LoRA 在线服务的管理成本。在产品和实验场景中，同一个基础模型之上可能同时存在大量任务版本、用户分支或策略变体。MinT 将可被请求的 LoRA 名称目录、驻留在主机内存中的 LoRA 参数，以及正在 GPU 上服务的 LoRA adapter 分开管理，从而避免所有版本都同时占用昂贵的 GPU 资源。通过这种分层管理方式，系统能够支持大规模 LoRA 目录下的精确版本选择，并进一步推导出百万级 LoRA 版本管理的容量模型。

% 第三，MinT 提供了端到端的 基于LoRA 的后训练工作流，包含训练更新、服务发布、版本管理和实验复现。系统将训练得到的 LoRA adapter 作为独立版本进行管理，负责完成参数导出、格式转换、基础模型兼容性检查、服务接入、rollout 记录和版本回滚。同时，MinT 提供 Tinker-compatible API，并通过 mint-cookbook recipes 覆盖 SFT、preference optimization、rollout-based RL 和 AutoResearch 等典型任务，使用户能够在统一接口下复现和迁移不同类型的后训练实验。
